\setcounter{section}{9}
\section{Graphs}

\begin{table}[h]
    \centering
    \caption{Graph Terminology}
    \begin{tabular}{llcc}
        \toprule
        \textbf{Type} & \textbf{Edges} & \textbf{Allow Multiple Edges?} & \textbf{Allow Loops?} \\
        \midrule
        Simple graph & Undirected & No & No \\
        Multigraph & Undirected & Yes & No \\
        Pseudograph & Undirected & Yes & Yes \\
        Simple directed graph & Directed & No & No \\
        Directed multigraph & Directed & Yes & Yes \\
        Mixed graph & Both & Yes & Yes \\
        \bottomrule
    \end{tabular}
\end{table}
Multiple edge mean multiple edge between the same pair of vertices. A loop is an edge that connects a vertex to itself.

\subsection{Basic knowledge}
The sum of the degrees of all vertices is 2 times the number of edges. Therefore the sum is always even.

This means the number of vertices with odd degree is always even.

The sum of in-degrees is equal to the sum of out-degrees, and both are equal to the number of edges.

\subsection{Common Graphs}
\plainBox{
    \subsubsection{Complete Graphs $K_n$}
    Graph with one edge between each pair of vertices.
    \begin{center}
        \textbf{Vertices:} $n$ \hfill
        \textbf{Edges:} $\frac{n(n-1)}{2}$ \hfill
        \textbf{Degrees of vertices:} $n-1$
    \end{center}
}

\plainBox{
    \subsubsection{Cycles $C_n$}
    Graph where the vertices are connected in a closed chain. Only defined for $n \geq 3$.
    \begin{center}
        \textbf{Vertices:} $n$ \hfill
        \textbf{Edges:} $n$ \hfill
        \textbf{Degrees of vertices:} $2$
    \end{center}
}

\plainBox{
    \subsubsection{Wheels $W_n$}
    A cycle with an additional center vertex connected to all other vertices. Only defined for $n \geq 4$.
    \begin{center}
        \textbf{Vertices:} $n$ \hfill
        \textbf{Edges:} $2(n-1)$ \hfill
        \textbf{Degrees of vertices:} Center: $n-1$, Other: $3$
    \end{center}
}

\plainBox{
    \subsubsection{n-cubes $Q_n$}
    A cube of dimension $n$. A cube of dimension $n$ can be made by connecting the corresponding vertices of two cubes of dimension $n-1$.
    \begin{center}
        \textbf{Vertices:} $2^n$ \hfill
        \textbf{Edges:} $n2^{n-1}$ \hfill
        \textbf{Degrees of vertices:} $n$
    \end{center}
}

